\documentclass[12pt]{article}
\setlength{\topmargin}{-0.5 in}
\setlength{\textheight}{9.5 in}
\setlength{\oddsidemargin}{-0.4 in}
\setlength{\textwidth}{7.0 in}
\usepackage{amsmath,graphicx,amssymb, diagbox}
\usepackage{tikz}
\usepackage{braket}
\newcommand*\circled[1]{\tikz[baseline=(char.base)]{
		\node[shape=circle,draw,inner sep=2pt] (char) {#1};}}
\pagestyle{empty}
\begin{document}
	
\centerline{QC HW 5
	\hspace{0.4 in}V. Kamat
	\hspace{0.4 in}Spring 2026}

\begin{enumerate}
\item Consider the \textbf{SWAP} unitary $U$ that, as the name implies, acts on any two arbitrary qubit states $|\psi\rangle$ and $|\phi\rangle$ as follows:
$$U\left(|\psi\rangle \otimes |\phi\rangle\right)=|\phi\rangle\otimes |\psi\rangle.$$
\begin{enumerate}
    \item Write a matrix for $U$ by evaluating its action on the four computational basis states.
    \item Prove that $U^2=I$. (This should make intuitive sense, given the definition.)
    \item Construct a quantum circuit implementing the swap unitary using \textbf{only} $\text{CNOT}$ gates. Verify (using bilinearity) that your quantum circuit correctly maps $|\psi\rangle|\phi\rangle\to |\phi\rangle|\psi\rangle$ for arbitrary qubit states $|\psi\rangle$ and $|\phi\rangle$. 
    \item The no-cloning theorem says that a cloning map $|\psi\rangle|0\rangle\to |\psi\rangle|\psi\rangle$ cannot be implemented by any unitary. Show that a weaker \textbf{state transfer} map $|\psi\rangle|0\rangle\to |0\rangle|\psi\rangle$ is achievable. 
\end{enumerate}
\item Recall that the Deutsch black box for a function $f:\{0,1\}\to \{0,1\}$ acts as
\[
O_f |x\rangle|y\rangle = |x\rangle|y \oplus f(x)\rangle.
\]


In the standard presentation of Deutsch's algorithm, the initial state is $|0\rangle|1\rangle$. In this problem, we investigate what happens if the second qubit is instead initialized to $|0\rangle$.

\begin{enumerate}
    \item[(a)] Start with initial state $|0\rangle|0\rangle$. Apply a Hadamard gate to both qubits and write the resulting state explicitly.
    \item[(b)] Apply $O_f$ to your state from part (a). 
    Express the result in terms of $f(0)$ and $f(1)$.
    \item[(c)] Apply a Hadamard gate to the first qubit, and write out the resulting state. 
    \item[(d)] If you measure the first qubit of the state from part (c), can you determine whether $f$ is constant or balanced with one query? Explain. 
\end{enumerate}
\item The one qubit gate $R_y(\theta)$ gate can be mathematically described as below:
$$R_y(\theta)=\begin{bmatrix}
\cos(\theta/2) &  -\sin (\theta/2)\\
\sin(\theta/2) & \cos (\theta/2)
\end{bmatrix}$$
Its name signifies that its action on a qubit amounts to a rotation on the Bloch sphere by an angle of $\theta$ around the $y$-axis. 
Prove that 
$$[R_y(\theta_A) \otimes R_y(\theta_B)]|\Phi^+\rangle=\dfrac{\cos\left(\frac{\theta_A-\theta_B}{2}\right)\left(|00\rangle +|11\rangle \right) + \sin\left(\frac{\theta_A-\theta_B}{2}\right) \left( |10\rangle - |01\rangle \right)}{\sqrt{2}},$$
where $|\Phi^+\rangle=\dfrac{1}{\sqrt{2}}\left(|00\rangle+|11\rangle\right)$ is the first Bell state. 
\end{enumerate}
\end{document}